\documentclass[a4paper, 11pt]{article}
\usepackage{xeCJK}
\usepackage{geometry}
\usepackage{titlesec}

\geometry{margin=1in}
\setCJKmainfont{MS Mincho}

% \title{進捗報告}
% \author{コレドール・パブロ}
% \date{２０２５年１０月１０日}

\titlespacing*{\section}{0pt}{2ex}{1ex}
\titlespacing*{\subsection}{0pt}{1ex}{0.8ex}
\titlespacing*{\subsubsection}{0pt}{1ex}{0.5ex}

\begin{document}

\begin{center}
    {\LARGE \bfseries 進捗報告} \\ % Title: Large and bold
    \vspace{0.3em} % A small vertical space
    Juan Pablo Corredor \\ % Author
    ２０２５年１０月１７日 % Date
\end{center}
\vspace{0.5em} % Adjust space between title block and main text

\section*{近況}
% Write your recent findings here.
\begin{itemize}
    \item スペインの「カタロニアシネマ・映像大学」の修士卒プロジェクトのBGMを制作。
    \item 藝大アニメーション科のプロジェクトのBGMも制作。
    \item 音響創造演習のプレゼンを完了。
    \item AES学生支部, NIME, ICMC等に論文の提出を準備。
    \item ドイツ語能力試験のC1のための勉強を開始。
    \item インターンシップの応募が続いている。
    \item 出身大学の後輩に、卒業研究のJUCEとC++について相談された。
\end{itemize} 

\section*{やったこと}
\begin{itemize}
    \item FFTウィンドウの作成方法を比較するための準備を完了。
    \item データをCSVファイルに書き込む方法を改善。
\end{itemize}


\section*{考えていること}

周波数の変更は期待通り音源のはじめに多く、時間が経てばたつほど、安定になる。

CSVファイルから音源を生成するプログラムはまだ完了ではないため、音源は来週になる。

% \begin{figure}
% \centering
% \includegraphics[width = 0.9\linewidth]{../../media/2025_02/EGuit_Harmonics_82Hz.png}
% \caption{８２Hzのエレキギターの倍音エンベロープ}
% \label{fig:eGuitHarmonics}
% \end{figure}

% \begin{figure}
% \centering
% \includegraphics[width = 0.9\linewidth]{../../media/2025_02/AGuit_Harmonics_146.png}
% \caption{１４６Hzのアコギターの倍音エンベロープ}
% \label{fig:aGuitHarmonics}
% \end{figure}

% \begin{figure}
% \centering
% \begin{subfigure}{1.0\textwidth}
%   \centering
%   \includegraphics[width=0.9\linewidth]{../media/training_data/20250613/elec_pitch.png}
%   \caption{トレーニングデータとして使用しているギターサンプル}
%   \label{fig:eGuitWave}
% \end{subfigure}%
% \vspace{5mm}
% \begin{subfigure}{1.0\textwidth}
%   \centering
%   \includegraphics[width=0.9\linewidth]{../media/training_data/20250613/elec_spectro.png}
%   \caption{そのスペクトルグラム}
%   \label{fig:eGuitSpectro}
% \end{subfigure}
% \caption{エレキギター}
% \label{fig:elecGuit}
% \end{figure}

\section*{やること}
% Outline future steps here.
\begin{itemize}
    \item 機械学習のモデルの開発に戻る。
    \item コロンビアの友人・知り合いを連絡、撥弦楽器の録音を依頼する。
\end{itemize}

% Add your bibliography here.
\bibliography{../../ref}
\bibliographystyle{apalike}

\end{document}